\documentclass[conference]{IEEEtran}

% =====================================================================
% Packages
% =====================================================================
\usepackage{amsmath,amssymb}
\usepackage{graphicx}
\usepackage{booktabs}
\usepackage{array}
\usepackage{tabularx}
\usepackage{xcolor}
\usepackage{tcolorbox}
\usepackage{algorithm}
\usepackage{algpseudocode}
\usepackage{microtype}
\usepackage{stfloats}
\usepackage{flushend}
\usepackage{cite}
\usepackage[hidelinks]{hyperref}

\tcbuselibrary{skins,breakable}
\graphicspath{{../figures/}{figures/}{scripts/figures/}}
\DeclareGraphicsExtensions{.pdf,.png}

% Search the sibling figures folder first; prefer PDF, with PNG fallback.
\newcommand{\paperfigure}[2]{%
    \includegraphics[width=#2]{#1}%
}

% Allow text to flow naturally around figures.
\raggedbottom

% Permit more figures on text pages.
\setcounter{topnumber}{3}
\setcounter{bottomnumber}{2}
\setcounter{totalnumber}{5}
\setcounter{dbltopnumber}{2}

\renewcommand{\topfraction}{0.90}
\renewcommand{\bottomfraction}{0.80}
\renewcommand{\textfraction}{0.08}
\renewcommand{\floatpagefraction}{0.70}
\renewcommand{\dbltopfraction}{0.90}
\renewcommand{\dblfloatpagefraction}{0.70}

% Compact, consistent float spacing.
\setlength{\textfloatsep}{10pt plus 2pt minus 2pt}
\setlength{\floatsep}{8pt plus 2pt minus 2pt}
\setlength{\intextsep}{8pt plus 2pt minus 2pt}
\setlength{\dbltextfloatsep}{10pt plus 2pt minus 2pt}
\setlength{\dblfloatsep}{8pt plus 2pt minus 2pt}

% =====================================================================
% Restrained colors and column-breakable highlight boxes
% =====================================================================
\definecolor{paperblue}{RGB}{35,75,110}
\definecolor{paperteal}{RGB}{20,105,90}
\definecolor{papergray}{RGB}{85,90,95}

\tcbset{
    enhanced,
    breakable,
    sharp corners,
    boxrule=0pt,
    leftrule=1.5pt,
    left=6pt,
    right=6pt,
    top=5pt,
    bottom=5pt,
    before skip=7pt,
    after skip=7pt,
    fonttitle=\bfseries\small,
    fontupper=\small
}

% Permit column breaks instead of leaving gaps before long boxes.
\newtcolorbox{methodbox}[1]{
    colback=paperblue!4,
    colframe=paperblue,
    coltitle=paperblue,
    colbacktitle=paperblue!4,
    title={#1},
    attach title to upper={\par\smallskip}
}

\newtcolorbox{resultbox}[1]{
    colback=paperteal!4,
    colframe=paperteal,
    coltitle=paperteal,
    colbacktitle=paperteal!4,
    title={#1},
    attach title to upper={\par\smallskip}
}

\newtcolorbox{scopebox}[1]{
    colback=black!3,
    colframe=papergray,
    coltitle=papergray,
    colbacktitle=black!3,
    title={#1},
    attach title to upper={\par\smallskip}
}

% =====================================================================
% Table columns and notation
% =====================================================================
\newcolumntype{Y}{>{\raggedright\arraybackslash}X}
\newcolumntype{Z}{>{\raggedleft\arraybackslash}X}

\newcommand{\R}{\mathbb{R}}
\newcommand{\trans}{^{\mathsf T}}
\newcommand{\norm}[1]{\left\lVert#1\right\rVert}

\DeclareMathOperator*{\argmin}{arg\,min}
\DeclareMathOperator{\col}{col}
\DeclareMathOperator{\tr}{tr}
\DeclareMathOperator{\range}{range}
\DeclareMathOperator{\clip}{clip}

% =====================================================================
% Title
% =====================================================================
\title{Control Density Functions for Whole-Arm
Obstacle Avoidance and Trajectory Tracking}

\author{
\IEEEauthorblockN{Author Names}
\IEEEauthorblockA{
Department and Institution\\
City, Country\\
Email addresses}
}

\begin{document}

\maketitle

\begin{abstract}
These notes present a control density function quadratic program
for whole-arm obstacle avoidance and Cartesian trajectory tracking.
A navigation density combines smooth obstacle factors with a
regularized target weighting. Parameterizing the density-weighted
control flux produces an affine transport constraint and a convex
quadratic program with joint-velocity bounds. The pointwise
optimization admits a reduction from $168$ flux coefficients
to at most seven coordinates. The formulation is evaluated on
a six-joint UF850 manipulator with arm and gripper collision
geometry. A tracking experiment reaches a final position error
of $4.84$\,mm while retaining positive recorded obstacle
clearance. The optimized implementation achieves a median
nominal-controller and filter time of $1.034$\,ms in a CPU
replay benchmark. Additional rollouts at $50$, $100$, and
$240$\,Hz retain positive clearance without measured
computation-budget overruns. These results establish finite-run
kinematic performance; physical robot execution and closed-loop
safety guarantees remain to be validated.
\end{abstract}

\begin{IEEEkeywords}
Control density functions, quadratic programming, robot
manipulators, obstacle avoidance, trajectory tracking.
\end{IEEEkeywords}

% =====================================================================
\section{Introduction}
\label{sec:introduction}
% =====================================================================

\begin{methodbox}{Main Contributions}
This work develops a density-based command filter for
robot manipulator motion in cluttered environments.
The main contributions are:
\begin{enumerate}
    \item \textbf{Finite-dimensional CDF-QP and computational
    reduction.}
    A quadratic weighted-flux parameterization yields
    one aggregate density transport constraint alongside
    joint-rate bounds. The pointwise problem is reduced
    from $168$ coefficients to at most seven coordinates
    without changing its optimum.

    \item \textbf{Whole-arm trajectory tracking.}
    Smooth factors for all robot-sphere--obstacle-sphere
    pairs are combined into one density. Arm and gripper
    geometry thereby influence the filtering of a nominal
    Cartesian pose-tracking command.

    \item \textbf{Simulation and timing evaluation.}
    UF850 experiments assess tracking, obstacle clearance,
    intervention, and numerical feasibility. Optimized CPU
    benchmarks and multiple update rates assess computation
    time relative to the available step budgets.
\end{enumerate}
\end{methodbox}

The central modeling choice is to combine all
robot--obstacle sphere pairs in one obstacle factor.
The resulting QP contains a single scalar density
transport inequality, rather than an inequality for
each pair. The number of pairs still affects geometry
and derivative evaluation.

The formulation is evaluated as a pointwise command
filter. The reported results demonstrate empirical
obstacle avoidance and tracking; they do not establish
sampled-data safety of the recomputed controller.

% =====================================================================
\section{Preliminaries and Density Construction}
\label{sec:preliminaries}
% =====================================================================

\subsection{Navigation Problem}

Consider the control-affine system
\begin{equation}
    \dot x=f(x)+g(x)u,
    \label{eq:system}
\end{equation}
where $x\in\mathcal X\subseteq\R^n$,
$u\in\mathcal U\subseteq\R^m$, and the vector fields
are continuously differentiable. Let
$\mathcal X_0$, $\mathcal X_T$, and $\mathcal X_u$
denote the initial, target, and unsafe sets.

The navigation objective is to reach $\mathcal X_T$
from $\mathcal X_0$ while avoiding $\mathcal X_u$
and respecting the input constraints. In trajectory
tracking, the objective is to follow a prescribed
reference while maintaining obstacle clearance.

Density-based navigation encodes target and unsafe
regions through a scalar field. The construction
is motivated by an occupancy interpretation:
desirable trajectories should avoid occupying
the unsafe region and approach the target
\cite{moyalan2024synthesizing,moyalan2024robust}.
A suitable condition on density transport is also
required; the scalar field alone is not a certificate.

\subsection{Smooth Obstacle Factors}

Let the unsafe set be
\begin{equation}
    \mathcal X_u=
    \bigcup_{\ell=1}^{N_c}
    \{x:d_\ell(x)\leq0\},
\end{equation}
where $N_c$ counts modeled collision conditions.
Positive signed clearance $d_\ell$ indicates
separation, zero indicates contact, and negative
clearance indicates penetration.

For a sensing width $\delta_\ell>0$, define
\begin{equation}
    S(s)=
    \begin{cases}
        0, &s\leq0,\\[1mm]
        \displaystyle
        \frac{\exp(-1/s)}
        {\exp(-1/s)+\exp[-1/(1-s)]},
        &0<s<1,\\[3mm]
        1, &s\geq1.
    \end{cases}
    \label{eq:smoothstep}
\end{equation}
The inverse obstacle bump is
\begin{equation}
    \Psi_\ell(x)=
    S\left(\frac{d_\ell(x)}{\delta_\ell}\right).
\end{equation}
It vanishes on the unsafe region, transitions through
the sensing region, and equals one beyond that region.

Combining the factors gives
\begin{equation}
    \Psi(x)=\prod_{\ell=1}^{N_c}\Psi_\ell(x).
    \label{eq:psi}
\end{equation}
Thus, $\Psi=0$ whenever any modeled collision condition
is active. Outside all sensing regions, $\Psi=1$.
For an obstacle-free scene, the product is defined
as one.

\subsection{Target Weighting}

Let $V(x)$ be positive away from the target.
A navigation density takes the form
\begin{equation}
    \rho(x)=\frac{\Psi(x)}{V(x)^\alpha},
    \qquad \alpha>0.
    \label{eq:density}
\end{equation}
The numerator encodes obstacle geometry and the
denominator weights the target region
\cite{moyalan2024synthesizing,moyalan2024robust}.

The robot implementation uses
\begin{equation}
    V_\epsilon(q)=
    \norm{p_{\mathrm{TCP}}(q)-p_g}^2+\epsilon,
    \qquad \epsilon>0,
    \label{eq:target}
\end{equation}
where $p_{\mathrm{TCP}}$ is the tool-center-point
position and $p_g$ is the final Cartesian goal.
The implemented density is
\begin{equation}
    \rho(q)=\frac{\Psi(q)}{V_\epsilon(q)^\alpha}.
    \label{eq:robot-density}
\end{equation}

Regularization keeps the density finite at the goal.
Theoretical results requiring a singular target
density therefore need additional justification
for this implementation.

% =====================================================================
\section{Robot Model and Nominal Controller}
\label{sec:robot}
% =====================================================================

\subsection{Euler--Lagrange Equations}

For a fully actuated robot with $n_q$ joints,
let $q\in\R^{n_q}$ denote joint position and
$v=\dot q$ joint velocity.
The kinetic energy and Lagrangian are
\begin{equation}
    \mathcal T(q,v)=\frac12v\trans M(q)v,
    \qquad L=\mathcal T-P(q).
\end{equation}
The Euler--Lagrange equations give
\begin{equation}
    \frac{\mathrm d}{\mathrm dt}
    \frac{\partial L}{\partial\dot q}
    -
    \frac{\partial L}{\partial q}
    =
    \tau-F(\dot q).
\end{equation}
Equivalently,
\begin{equation}
    M(q)\ddot q+C(q,\dot q)\dot q
    +G(q)+F(\dot q)=\tau,
    \label{eq:el}
\end{equation}
where $M(q)=M(q)\trans\succ0$ is inertia,
$C(q,\dot q)\dot q$ contains Coriolis and centrifugal
terms, $G(q)=\nabla_qP(q)$ is gravity, and $F$
represents friction.

With
\begin{equation}
    x=\begin{bmatrix}q\\v\end{bmatrix},
    \qquad h(q,v)=C(q,v)v+G(q)+F(v),
\end{equation}
the control-affine representation is
\begin{equation}
    \dot x=
    \underbrace{
    \begin{bmatrix}
        v\\-M(q)^{-1}h(q,v)
    \end{bmatrix}}_{f(x)}
    +
    \underbrace{
    \begin{bmatrix}
        0\\M(q)^{-1}
    \end{bmatrix}}_{g(x)}
    \tau.
    \label{eq:el-affine}
\end{equation}

\subsection{Velocity-Level Implementation}

The evaluated controller uses
\begin{equation}
    \dot q=u,\qquad q,u\in\R^6,
    \label{eq:kinematic}
\end{equation}
with componentwise bounds
\begin{equation}
    -u_{\max}\mathbf1\leq u\leq u_{\max}\mathbf1.
    \label{eq:bounds}
\end{equation}
Here $u$ is a joint-velocity command.

The simulation integrates
\begin{equation}
    q_{k+1}=q_k+\Delta t\,u_k
    \label{eq:integration}
\end{equation}
and assigns the resulting joint positions directly
to the PyBullet robot. A physical implementation
would require an inner velocity servo that tracks
the command accurately.

\begin{scopebox}{Evaluated model}
The Euler--Lagrange equations describe the physical
arm. The experiments evaluate joint-velocity
integration, collision geometry, and computation
time. Motor dynamics and hardware tracking are
not evaluated.
\end{scopebox}

\subsection{Nominal Pose-Tracking Controller}
\label{sec:nominal}

The nominal controller is OSCBF's
\texttt{PoseTask}\allowbreak\texttt{Velocity}\allowbreak\texttt{Controller}~\cite{morton2025oscbf}.
It tracks TCP position and orientation without
considering obstacles. The CDF-QP subsequently
filters its joint-velocity command.

\subsubsection{Reference}

The TCP follows a straight Cartesian segment
with minimum-jerk timing:
\begin{equation}
    p_{\mathrm{ref}}(t)
    =p_0+s(\sigma)(p_g-p_0),
\end{equation}
where $p_0=p_{\mathrm{TCP}}(q_0)$,
$\sigma=\clip(t/T_f,0,1)$, and
\begin{equation}
    s(\sigma)=10\sigma^3-15\sigma^4+6\sigma^5.
\end{equation}
For $0\leq t\leq T_f$,
\begin{equation}
    \dot p_{\mathrm{ref}}(t)=
    \frac{30\sigma^2-60\sigma^3+30\sigma^4}{T_f}
    (p_g-p_0).
\end{equation}
The duration is $T_f=7$\,s. Reference velocity
and acceleration vanish at both endpoints.
The desired tool orientation is fixed:
\begin{equation}
    R_d=R(q_0).
\end{equation}

\subsubsection{Position and Orientation Feedback}

The linear velocity command is
\begin{equation}
    v_{\mathrm{cmd}}
    =
    \dot p_{\mathrm{ref}}
    +K_p\bigl(p_{\mathrm{ref}}-p_{\mathrm{TCP}}(q)\bigr).
\end{equation}

Let $R=[r_1,r_2,r_3]$ and
$R_d=[r_{d,1},r_{d,2},r_{d,3}]$ map tool-frame
coordinates to the world frame.
The controller's orientation-error vector is
\begin{equation}
    \mathbf e_R
    =
    -\frac12\sum_{i=1}^{3}
    r_i\times r_{d,i}.
    \label{eq:orientation-vector}
\end{equation}
This vector is expressed in the world frame,
consistently with the angular Jacobian.
With zero desired angular velocity,
\begin{equation}
    \omega_{\mathrm{cmd}}=-K_R\mathbf e_R.
\end{equation}
The gains are $K_p=1.4$ and $K_R=1.0$.

\subsubsection{Joint-Velocity Mapping}

Let $J(q)\in\R^{6\times6}$ be the world-frame
TCP geometric Jacobian, ordered as linear velocity
followed by angular velocity.
The mass-weighted generalized inverse is
\begin{equation}
    J^\#=
    M^{-1}J\trans
    \left(JM^{-1}J\trans\right)^{-1}.
    \label{eq:jacobian-inverse}
\end{equation}
The nominal command is
\begin{equation}
\begin{aligned}
    \widehat u_{\mathrm{nom}}
    &=
    J^\#
    \begin{bmatrix}
        v_{\mathrm{cmd}}\\
        \omega_{\mathrm{cmd}}
    \end{bmatrix},\\
    u_{\mathrm{nom}}
    &=
    \clip\left(
        \widehat u_{\mathrm{nom}},
        -u_{\max}\mathbf1,
        u_{\max}\mathbf1
    \right).
    \label{eq:nominal-command}
\end{aligned}
\end{equation}
Clipping is componentwise, with
$u_{\max}=0.7$\,rad/s.

The mass-weighted branch uses the direct inverse
in \eqref{eq:jacobian-inverse}, without a damping
term. It therefore requires a nonsingular task
inertia matrix; clipping alone does not address
singularity handling.

The controller applies a null-space posture task
only when the number of joints exceeds the
six-dimensional pose-task dimension. Consequently,
the configured posture gain contributes no additional
command for this six-joint experiment.

\begin{methodbox}{Roles of the two controllers}
The nominal controller supplies Cartesian tracking
and orientation regulation. The CDF-QP minimizes
deviation from that command while imposing the
pointwise density transport condition and joint-rate
bounds. Tracking is an objective, rather than
a separate hard constraint.
\end{methodbox}

% =====================================================================
\section{CDF-QP and Finite-Dimensional Approximation}
\label{sec:cdf}
% =====================================================================

\subsection{Density Transport}

For a differentiable feedback field, define
\begin{equation}
    R(x)=
    \nabla_x\cdot
    \left(\rho(x)[f(x)+g(x)u(x)]\right).
    \label{eq:residual}
\end{equation}
CDF-based navigation conditions constrain this
residual over an operating domain under the
assumptions of the associated theorem
\cite{moyalan2024synthesizing,moyalan2024robust}.
A positive numerical margin is written as
\begin{equation}
    R(x)\geq\eta,\qquad \eta>0.
\end{equation}

The divergence contains spatial derivatives
of the feedback, motivating a parameterized
feedback field.
Introduce
\begin{equation}
    w(x)=\rho(x)u(x),\qquad
    d_\rho(x)=\nabla_x\cdot(\rho(x)f(x)).
\end{equation}
With $g=[g_1,\ldots,g_m]$,
\begin{equation}
\begin{aligned}
    R(x)={}&d_\rho(x)\\
    &+\sum_{j=1}^{m}
    \left[
        g_j(x)\trans\nabla_xw_j(x)
        +w_j(x)\nabla_x\cdot g_j(x)
    \right].
    \label{eq:flux-residual}
\end{aligned}
\end{equation}

For the Euler--Lagrange model,
$\nabla_x\cdot g_j=0$, since $M$ depends only
on $q$. With torque-weighted flux $w=\rho\tau$,
\begin{equation}
    R(q,v)=d_\rho(q,v)
    +\sum_{j=1}^{n_q}
    e_j\trans M(q)^{-1}\nabla_vw_j(q,v).
\end{equation}
This identity describes the torque-level formulation.

For the implemented velocity model,
$f=0$ and $g=I_6$, giving
\begin{equation}
    R(q)=\nabla_q\cdot w(q).
\end{equation}

\subsection{Quadratic Flux Basis}

Set $z=q/\pi$ and collect monomials through
total degree two:
\begin{equation}
    \Phi_2(q)=
    \col\left(
        1,\{z_i\}_{i=1}^{6},
        \{z_i z_j\}_{1\leq i\leq j\leq6}
    \right)
    \in\R^{28}.
\end{equation}
There are one constant, six linear,
and twenty-one quadratic terms.

Define
\begin{equation}
    b(q)=\Psi(q)\Phi_2(q),\qquad
    w_j(q;\theta_j)=b(q)\trans\theta_j,
\end{equation}
where $\theta_j\in\R^{28}$.
Stacking all channels gives
\begin{equation}
    \theta=\col(\theta_1,\ldots,\theta_6)\in\R^{168}.
\end{equation}

Where $\rho>0$, the recovered input is
\begin{equation}
    u_j(q;\theta_j)
    =
    \frac{w_j(q;\theta_j)}{\rho(q)}
    =
    V_\epsilon(q)^\alpha
    \Phi_2(q)\trans\theta_j.
\end{equation}
The obstacle factor cancels in the quotient,
but remains in the transport constraint and
influences the optimized coefficients.

\subsection{Command and Transport Maps}

For numerical evaluation, define
\begin{equation}
    \bar\rho(q)=\max\{\rho(q),10^{-10}\},
\end{equation}
and
\begin{equation}
    B(q)=I_6\otimes
    \frac{b(q)\trans}{\bar\rho(q)}
    \in\R^{6\times168}.
\end{equation}
The command is $u_\theta=B\theta$.

Define the transport row
\begin{equation}
    a(q)\trans=
    \left[
        \left(\frac{\partial b}{\partial q_1}\right)\trans
        \ \cdots\
        \left(\frac{\partial b}{\partial q_6}\right)\trans
    \right].
\end{equation}
Then
\begin{equation}
    R_\theta(q)
    =
    \sum_{j=1}^{6}\frac{\partial w_j}{\partial q_j}
    =
    a(q)\trans\theta.
    \label{eq:linear-residual}
\end{equation}
All derivatives hold $\theta$ fixed.

\subsection{Finite-Dimensional QP}

At the current configuration, solve
\begin{equation}
\begin{aligned}
    \theta^\star=\argmin_\theta\quad&
    \frac12\norm{B\theta-u_{\mathrm{nom}}}^2\\
    &+\frac{\lambda_u}{2}\norm{B\theta}^2
    +\frac{\lambda_\theta}{2}\norm{\theta}^2\\
    \text{subject to}\quad&
    a\trans\theta\geq\eta,\\
    &-u_{\max}\mathbf1
    \leq B\theta\leq u_{\max}\mathbf1.
    \label{eq:qp}
\end{aligned}
\end{equation}
The filtered command is $u^\star=B\theta^\star$.

The objective preserves the nominal command,
penalizes command magnitude, and regularizes
the coefficients. Its Hessian and linear term are
\begin{align}
    H&=(1+\lambda_u)B\trans B+\lambda_\theta I,\\
    c&=-B\trans u_{\mathrm{nom}}.
\end{align}
For $\lambda_\theta>0$, $H$ is positive definite,
so a feasible QP has a unique optimizer.

There are $168$ coefficients and $13$ scalar
inequalities: one aggregate transport condition
and twelve joint-rate bounds.

\subsection{Exact Pointwise Reduction}

The command and constraint depend on $\theta$
through six input values and one residual.
Define
\begin{equation}
    A(q)=
    \begin{bmatrix}
        B(q)\\a(q)\trans
    \end{bmatrix}
    \in\R^{7\times168}.
\end{equation}
Let $Q\in\R^{168\times r}$ have orthonormal
columns spanning $\range(A\trans)$, with $r\leq7$.
Any coefficient vector decomposes as
\begin{equation}
    \theta=Q\xi+\theta_\perp,\qquad
    A\theta_\perp=0.
\end{equation}
The component $\theta_\perp$ changes neither
the command nor the residual, but increases
$\norm{\theta}^2$. Therefore,
$\theta_\perp=0$ at the optimum.

Setting
\begin{equation}
    \widetilde B=BQ,\qquad
    \widetilde a=Q\trans a
\end{equation}
gives
\begin{equation}
\begin{aligned}
    \xi^\star=\argmin_{\xi\in\R^r}\quad&
    \frac12\norm{\widetilde B\xi-u_{\mathrm{nom}}}^2\\
    &+\frac{\lambda_u}{2}\norm{\widetilde B\xi}^2
    +\frac{\lambda_\theta}{2}\norm{\xi}^2\\
    \text{subject to}\quad&
    \widetilde a\trans\xi\geq\eta,\\
    &-u_{\max}\mathbf1
    \leq\widetilde B\xi
    \leq u_{\max}\mathbf1.
    \label{eq:reduced}
\end{aligned}
\end{equation}
This reduction is exact for the stated pointwise
objective and constraints. It is performed after
evaluating the frozen-coefficient transport map.

The optimized implementation uses a seven-coordinate
QP and fused JAX evaluation. On the tested replay
samples, its maximum command difference from
the full formulation is
$1.11\times10^{-5}$\,rad/s, below the specified
$10^{-4}$\,rad/s tolerance.

\begin{methodbox}{Computational structure}
The flux field has $168$ coefficients, but the
pointwise optimization needs only six command values
and one transport residual. Coefficient directions
that affect neither quantity can be removed,
leaving at most seven optimization coordinates.
\end{methodbox}

\subsection{Scope of the Pointwise Condition}

The residual in \eqref{eq:linear-residual}
is evaluated with frozen coefficients.
The online controller recomputes
$\theta^\star(q,t)$, whose spatial derivatives
are absent from that expression.
The constraint is also checked only at sampled states.

The measured residual therefore does not establish
a domain-wide inequality or sampled-data invariance
for the recomputed policy.
The experiments report empirical tracking,
clearance, and numerical feasibility.

The QP has no slack and no general feasibility
guarantee. The density floor changes $w=\rho u$
wherever it is active. Self-collision and
joint-position constraints are not included.

% =====================================================================
\section{UF850 Setup and Implementation}
\label{sec:setup}
% =====================================================================

\subsection{Whole-Arm Geometry}

Robot sphere $a$ has center $p_a(q)$ and radius
$r_a$. Obstacle $b$ has center $o_b$ and radius
$R_b$. Their surface clearance is
\begin{equation}
    d_{ab}(q)=
    \norm{p_a(q)-o_b}-r_a-R_b.
\end{equation}
The obstacle factor is
\begin{equation}
    \Psi(q)=
    \prod_{a=1}^{N_r}\prod_{b=1}^{N_o}
    S\left(\frac{d_{ab}(q)}{\delta}\right).
    \label{eq:robot-psi}
\end{equation}

The closed-gripper model contains $20$ collision
spheres: $13$ arm/wrist spheres and seven additional
gripper-cover spheres attached to the wrist frame.
With five obstacle spheres, $100$ pairwise factors
enter \eqref{eq:robot-psi}, while the QP retains
one aggregate transport inequality.

Robot sphere centers are transformed from their
local link coordinates using differentiable
forward kinematics. The gripper cover is specific
to the fixed closed configuration.
The collision-model implementation provides the
local sphere centers and radii
\cite{xarmimplementation}.

Independent PyBullet mesh-distance and obstacle-contact
queries provide a second geometric measurement.
Sphere and mesh clearances need not agree because
they use different robot representations.

\subsection{Scene and Parameters}

The experiment uses a seven-second Cartesian reference
through five static obstacles.
The initial joint configuration is
\begin{equation}
    q_0=[0,-0.7,-1.0,0,0.8,0]\trans\ \mathrm{rad},
\end{equation}
and the final goal is
\begin{equation}
    p_g=[0.38,0.35,0.43]\trans\ \mathrm{m}.
\end{equation}
The gripper remains closed, and the TCP is
offset $0.172$\,m along the wrist tool axis.

Table~\ref{tab:setup} summarizes the setup.
Obstacle centers and radii are listed in
Appendix~\ref{app:obstacles}.

\begin{table}[t]
    \caption{Robot, scene, and controller setup.}
    \label{tab:setup}
    \centering
    \small
    \begin{tabularx}{\columnwidth}{@{}YZ@{}}
        \toprule
        Setting & Value\\
        \midrule
        Robot & UF850, six joints\\
        Gripper & Fixed closed configuration\\
        Robot collision spheres & $20$\\
        Obstacle spheres & $5$\\
        Robot--obstacle pairs & $100$\\
        Aggregate transport inequalities & $1$\\
        Reference duration $T_f$ & $7$ s\\
        TCP offset & $0.172$ m\\
        Position/orientation gains & $1.4/1.0$\\
        Joint-rate bound $u_{\max}$ & $0.7$ rad/s\\
        Sensing width $\delta$ & $0.10$ m\\
        Density exponent $\alpha$ & $0.2$\\
        Target regularization $\epsilon$ & $10^{-3}$\\
        Transport margin $\eta$ & $0.01$\\
        Input penalty $\lambda_u$ & $0.02$\\
        Coefficient penalty $\lambda_\theta$ & $2\times10^{-4}$\\
        Optimized rollout rates & $50$, $100$, $240$ Hz\\
        \bottomrule
    \end{tabularx}
\end{table}

The moving reference enters through
$u_{\mathrm{nom}}$. The density denominator
retains the fixed final goal.
Thus, this is a static-density filter with a
time-varying nominal command, rather than the
time-varying density construction in
\cite{narayanan2024dynamic}.

\begin{algorithm}[t]
\caption{Online Density-Based Joint-Velocity Filter}
\label{alg:cdf}
\begin{algorithmic}[1]
\Require Kinematics, collision geometry, reference,
QP parameters, and timestep $\Delta t$
\State Initialize geometry and the cached solver
\State Compile and warm up differentiable functions
\For{each sample $k$}
    \State Obtain $q_k$ and reference values at $t_k$
    \State Compute $u_{\mathrm{nom}}(q_k,t_k)$
    \State Evaluate clearances, $\Psi$, $\rho$, and basis derivatives
    \State Assemble command and transport maps
    \State Solve \eqref{eq:qp} or its reduction \eqref{eq:reduced}
    \If{the solve is unsuccessful}
        \State Flag failure and terminate the experiment
    \EndIf
    \State Recover the filtered command $u_k$
    \State Integrate $q_{k+1}=q_k+\Delta t\,u_k$
    \State Update simulated joint positions
    \State Record tracking, clearance, residual, and timing
\EndFor
\end{algorithmic}
\end{algorithm}

% =====================================================================
\section{Simulation Results}
\label{sec:results}
% =====================================================================

\subsection{Evaluation Protocol}

Two related evaluations are reported:
\begin{enumerate}
    \item \textbf{Tracking experiment:}
    a seven-second trajectory sampled at $240$\,Hz
    with $1681$ state samples and $1680$ commands.

    \item \textbf{Optimized implementation benchmark:}
    CPU replay timing and additional rollouts
    at $50$, $100$, and $240$\,Hz.
\end{enumerate}
These correspond to repository validations~9
and~10, respectively.

The coordinate, error, command, and intervention plots
describe the tracking experiment. The robot snapshots
show the optimized $100$\,Hz rollout.
Timing plots describe the optimized implementation.

\begin{figure*}[!t]
    \centering
    \paperfigure{fig08_robot_snapshots}{\textwidth}
    \caption{UF850 motion in the optimized $100$\,Hz
    rollout at $t=0$, $2.33$, $4.67$, and $7.00$\,s.
    The dashed curve is the TCP reference, the blue curve
    is the accumulated TCP path, and the green marker
    denotes the goal. The gripper remains closed.
    Snapshots are rendered from saved joint states.}
    \label{fig:snapshots}
\end{figure*}

\subsection{Robot Motion Through Clutter}

Figure~\ref{fig:snapshots} shows the robot raising
and reconfiguring its arm while approaching the goal.
The TCP path departs from the straight reference
during the middle of the motion and approaches
the goal in the final snapshot.

The optimized $100$\,Hz rollout ends with a
$4.76$\,mm position error. Its minimum sphere
and mesh clearances are $20.41$\,mm and
$48.49$\,mm, with zero recorded obstacle-contact
steps. These measurements complement the snapshots,
which show only four selected instants.

\subsection{Cartesian Tracking}

Figure~\ref{fig:coordinates} compares the TCP
coordinates with the moving reference in the
tracking experiment.
The largest visible deviation occurs in $y$
during the middle of the motion.
The $z$ coordinate initially remains below its
reference and subsequently rises toward the
final height.

\begin{figure}[!tbp]
    \centering
    \paperfigure{fig01_tcp_tracking}{\columnwidth}
    \caption{TCP coordinates in the tracking experiment.
    Dashed curves denote the reference and solid curves
    denote the CDF-QP response. Time is simulated time.}
    \label{fig:coordinates}
\end{figure}

Position error is
\begin{equation}
    e_p(t)=
    \norm{p_{\mathrm{TCP}}(q(t))-p_{\mathrm{ref}}(t)}.
\end{equation}
The reported orientation error is the rotation angle
\begin{equation}
    e_{\mathrm{ang}}(t)=
    \cos^{-1}\left[
        \clip\left(
            \frac{\tr(R_d\trans R(q(t)))-1}{2},
            -1,1
        \right)
    \right].
\end{equation}
This scalar angle differs from the feedback vector
$\mathbf e_R$ in \eqref{eq:orientation-vector}.

Figure~\ref{fig:errors} shows errors increasing
during the middle of the motion and decreasing
near the final reference.
The position RMS and maximum errors are
$66.82$\,mm and $115.37$\,mm.
The orientation RMS and maximum errors are
$8.80^\circ$ and $16.06^\circ$.

The final errors decrease to $4.84$\,mm and
$1.59^\circ$. Thus, the experiment demonstrates
small terminal errors after a substantial
transient tracking deviation.

\begin{figure}[!tbp]
    \centering
    \paperfigure{fig02_tracking_errors}{\columnwidth}
    \caption{Position error relative to the moving
    reference and orientation angle relative to the
    fixed desired tool orientation. Time is simulated time.}
    \label{fig:errors}
\end{figure}

\begin{table}[!tbp]
    \caption{Tracking errors. RMS values use all
    $1681$ uniformly spaced state samples.}
    \label{tab:tracking}
    \centering
    \small
    \setlength{\tabcolsep}{4pt}
    \begin{tabularx}{\columnwidth}{@{}Yrrr@{}}
        \toprule
        Error & RMS & Maximum & Final\\
        \midrule
        Position (mm) & 66.82 & 115.37 & 4.84\\
        Orientation (deg) & 8.80 & 16.06 & 1.59\\
        \bottomrule
    \end{tabularx}
\end{table}

\subsection{Joint Commands and Intervention}

Figure~\ref{fig:joints} compares nominal and
filtered joint velocities.
The filter modifies several joints simultaneously,
particularly joints 1, 2, 3, and 5 during
the middle of the motion.
The maximum absolute command is $0.684$\,rad/s,
below the $0.7$\,rad/s bound.

\begin{figure*}[!t]
    \centering
    \paperfigure{fig03_joint_velocities}{\textwidth}
    \caption{Nominal and filtered joint velocities.
    Dashed curves denote nominal commands, solid curves
    denote filtered commands, and dotted lines mark
    the $\pm0.7$\,rad/s bounds.
    A command at $t_k$ acts over the following
    integration interval.}
    \label{fig:joints}
\end{figure*}

Define intervention as
\begin{equation}
    I_k=\norm{u_k-u_{\mathrm{nom},k}}_2.
\end{equation}
Its RMS and maximum are $0.178$\,rad/s
and $0.328$\,rad/s. The zero-order-hold integral is
\begin{equation}
    \Delta t\sum_k I_k=0.876\ \mathrm{rad}.
\end{equation}

The lower panel of Fig.~\ref{fig:safety}
shows intervention concentrated in the middle
of the motion, overlapping the tracking-error
transient.
Intervention includes the effects of regularization
and the transport constraint; it is not a
safety-constraint activation measure or a causal
comparison with an unfiltered rollout.

\subsection{Clearance and Transport Residual}

Both recorded clearance measures remain positive
in Fig.~\ref{fig:safety}.
Their minima are $20.37$\,mm for the collision
spheres and $48.54$\,mm for the PyBullet mesh queries.
No obstacle-contact steps are reported.

The residual remains near $\eta=0.01$ over
much of the trajectory and rises near the end.
No margin violations are recorded using a
$10^{-8}$ tolerance, and all QPs are reported
feasible.

\begin{figure}[!tbp]
    \centering
    \paperfigure{fig04_safety_intervention}{\columnwidth}
    \caption{Clearance, transport residual, and intervention.
    Top: minimum sphere and mesh clearances.
    Middle: raw CDF residual and margin $\eta=0.01$.
    Bottom: $\|u-u_{\mathrm{nom}}\|_2$.
    Clearance is evaluated after integration;
    residual and intervention are evaluated at the
    command state before integration.}
    \label{fig:safety}
\end{figure}

\begin{table}[t]
    \caption{Clearance, feasibility, and command modification
    in the tracking experiment.}
    \label{tab:clearance}
    \centering
    \small
    \begin{tabularx}{\columnwidth}{@{}YZ@{}}
        \toprule
        Metric & Value\\
        \midrule
        Minimum sphere clearance & $20.37$ mm\\
        Minimum mesh clearance & $48.54$ mm\\
        Obstacle-contact steps & $0$\\
        Minimum transport residual & $0.0100$\\
        Maximum joint speed & $0.684$ rad/s\\
        RMS intervention & $0.178$ rad/s\\
        Maximum intervention & $0.328$ rad/s\\
        Integrated intervention & $0.876$ rad\\
        All QPs feasible & Yes\\
        \bottomrule
    \end{tabularx}
\end{table}

\begin{resultbox}{Tracking and clearance}
The tracking experiment reaches a final position
error of $4.84$\,mm while retaining positive
recorded clearances and zero obstacle contacts.
The larger RMS and maximum errors quantify
the transient departure from the reference.
\end{resultbox}

% Let tracking floats and timing text share available column space.

\subsection{Optimized CPU Timing}

% Declare the wide figure before the detailed timing discussion.
\begin{figure*}[!t]
    \centering
    \paperfigure{fig06_optimized_timing}{\textwidth}
    \caption{Optimized CPU replay timing over $480$
    evaluations after warm-up.
    Left: empirical cumulative distributions on a
    logarithmic time axis.
    Right: median, $95$th-percentile, and maximum
    durations. Timing scopes are nested, not additive.
    Bars summarize measured samples rather than
    confidence intervals across independent experiments.}
    \label{fig:optimized-timing}
\end{figure*}

The optimized implementation was benchmarked on an
Intel Core i7-10700K CPU at $3.80$\,GHz using
JAX~0.4.30 and Clarabel~0.11.1, with one requested
BLAS thread.
The benchmark evaluates $160$ saved states
over three repeats, giving $480$ timed
evaluations after explicit warm-up.

Figure~\ref{fig:optimized-timing} and
Table~\ref{tab:timing} summarize the results.
The median nominal-controller and filter time is
$1.034$\,ms, the $95$th percentile is
$1.091$\,ms, and the maximum is $1.562$\,ms.
The filter median is $0.757$\,ms, while the
solver-call median is $0.057$\,ms.

\begin{table}[t]
    \caption{Optimized CPU replay timing in milliseconds.}
    \label{tab:timing}
    \centering
    \small
    \setlength{\tabcolsep}{4pt}
    \begin{tabularx}{\columnwidth}{@{}Yrrr@{}}
        \toprule
        Component & Median & p95 & Maximum\\
        \midrule
        Nominal + filter & 1.034 & 1.091 & 1.562\\
        Density filter & 0.757 & 0.811 & 1.287\\
        Solver only & 0.057 & 0.094 & 0.108\\
        \bottomrule
    \end{tabularx}
\end{table}

These timing scopes are nested: the filter includes
the solver, and nominal plus filter includes the
filter. They must not be added together.

Nominal-plus-filter timing includes reference
generation, nominal control, JAX evaluation and
host synchronization, matrix construction,
checks, QP update and solve, and filter-return
diagnostics.
It excludes Bullet stepping, rendering, sensor
input, and communication.

\subsection{Fixed-Rate Rollouts}

% Declare the wide figure before the detailed rollout discussion.
\begin{figure*}[!t]
    \centering
    \paperfigure{fig07_fixed_rate_budget}{\textwidth}
    \caption{Measured computation time as a percentage
    of the simulated step budget.
    Left: nominal controller and filter.
    Right: the loop including Bullet stepping and
    diagnostics. All measured maxima remain below
    the dashed $100\%$ budget line.
    The rollouts are headless and unpaced.}
    \label{fig:budget}
\end{figure*}

Additional rollouts use simulated update rates
of $50$, $100$, and $240$\,Hz, corresponding
to step budgets of $20$, $10$, and $4.167$\,ms.
Each follows the seven-second reference.

Figure~\ref{fig:budget} normalizes computation
time by the associated budget.
The controller timing includes nominal control
and the filter.
The diagnostic-loop timing additionally includes
joint integration, Bullet stepping, mesh and
contact queries, and logged diagnostics.

At $100$\,Hz, the maximum controller and loop
times are $1.409$\,ms and $2.245$\,ms.
At $240$\,Hz, they are $1.873$\,ms and
$3.373$\,ms, approximately $45\%$ and $81\%$
of the $4.167$\,ms budget.
No measured controller or diagnostic-loop
budget overruns occur in these runs.

\begin{table}[t]
    \caption{Fixed-rate computation times.
    Budgets and durations are in milliseconds.}
    \label{tab:rates}
    \centering
    \small
    \setlength{\tabcolsep}{3pt}
    \begin{tabularx}{\columnwidth}{@{}Yrrrr@{}}
        \toprule
        Rate & Budget & Ctrl.\ p95 & Ctrl.\ max & Loop max\\
        \midrule
        50 Hz  & 20.000 & 1.083 & 1.405 & 2.250\\
        100 Hz & 10.000 & 1.096 & 1.409 & 2.245\\
        240 Hz & 4.167  & 1.131 & 1.873 & 3.373\\
        \bottomrule
    \end{tabularx}
\end{table}

All runs reach the target, all QPs solve successfully,
and no obstacle contacts or density-floor
activations are recorded.
Minimum clearances and terminal errors remain
similar across update rates
(Table~\ref{tab:rate-safety}).
The maximum joint speed remains below
$0.7$\,rad/s at every tested rate.

\begin{table}[t]
    \caption{Minimum clearance and terminal error
    across update rates. All distances are in millimeters.}
    \label{tab:rate-safety}
    \centering
    \small
    \setlength{\tabcolsep}{4pt}
    \begin{tabularx}{\columnwidth}{@{}Yrrr@{}}
        \toprule
        Rate & Sphere & Mesh & Final error\\
        \midrule
        50 Hz  & 20.49 & 48.40 & 4.62\\
        100 Hz & 20.41 & 48.49 & 4.76\\
        240 Hz & 20.37 & 48.54 & 4.84\\
        \bottomrule
    \end{tabularx}
\end{table}

\begin{resultbox}{Optimized implementation}
The median nominal-plus-filter replay time is
$1.034$\,ms. All tested update rates retain positive
clearance without measured budget overruns.
The $100$\,Hz run uses at most $2.245$\,ms
of its $10$\,ms budget when diagnostics are included.
\end{resultbox}

\subsection{Interpretation and Remaining Validation}

The figures provide complementary evidence:
snapshots show robot configuration, tracking plots
quantify reference deviation, command plots show
joint-rate limits, and clearance queries assess
the modeled collision geometry.
Timing experiments assess computation separately
from tracking accuracy.

The rollouts are headless, unpaced kinematic
experiments. Measured durations below simulated
step budgets demonstrate computational headroom
for these finite runs, not hard real-time bounds.
Rendering, sensing, communication, scheduling
jitter, and physical motor tracking remain
outside the reported controller timing.

The experiments also use one scene and reference.
They do not establish robustness across initial
conditions or obstacle placements, general
sphere coverage, self-collision avoidance,
or superiority over a matched CBF baseline.

% Avoid forcing partially empty columns at this section boundary.

% The flushend package balances the final page automatically.
% =====================================================================
\section{Conclusion}
\label{sec:conclusion}
% =====================================================================

A density-based command filter was formulated
using smooth obstacle factors and a quadratic
weighted-flux basis.
The QP combines all modeled robot--obstacle pairs
in one aggregate transport condition alongside
joint-rate bounds.
Its $168$-coefficient pointwise formulation
admits reduction to at most seven coordinates.

UF850 simulations demonstrate tracking through
clutter with positive recorded clearance and
zero obstacle contacts.
The optimized implementation achieves
millisecond-scale CPU updates and remains
within the measured step budgets at $50$,
$100$, and $240$\,Hz.
Further work is required to analyze the
recomputed sampled controller and validate
geometry coverage and physical execution.

% =====================================================================
\appendices
\section{Obstacle Configuration}
\label{app:obstacles}
% =====================================================================

Table~\ref{tab:obstacles} lists the five static
obstacle spheres in the world frame.
Robot collision-sphere centers and radii are
specified by the implementation's
\texttt{make\_uf850\_collision\_model}
function \cite{xarmimplementation}.

\begin{table}[!ht]
    \caption{World-frame obstacle centers and radii.
    All values are in meters.}
    \label{tab:obstacles}
    \centering
    \small
    \setlength{\tabcolsep}{4pt}
    \begin{tabularx}{\columnwidth}{@{}Yrrrr@{}}
        \toprule
        Obstacle & $x$ & $y$ & $z$ & Radius\\
        \midrule
        1 & 0.5665 & 0.2553 & 0.3056 & 0.075\\
        2 & 0.4900 & 0.2844 & 0.5415 & 0.080\\
        3 & 0.2924 & 0.2063 & 0.2649 & 0.076\\
        4 & 0.6790 & 0.0892 & 0.4238 & 0.084\\
        5 & 0.5021 & 0.3967 & 0.3615 & 0.061\\
        \bottomrule
    \end{tabularx}
\end{table}

% =====================================================================
\begin{thebibliography}{9}
% =====================================================================

\bibitem{moyalan2024synthesizing}
J. Moyalan, S. S. K. S. Narayanan, A. Zheng, and U. Vaidya,
``Synthesizing controller for safe navigation using control
density function,''
arXiv:2403.14464, 2024.

\bibitem{moyalan2024robust}
J. Moyalan, S. S. K. S. Narayanan, and U. Vaidya,
``Control density function for robust safety and convergence,''
arXiv:2407.05133, 2024.

\bibitem{narayanan2024dynamic}
S. S. K. S. Narayanan, J. Moyalan, and U. Vaidya,
``Safe navigation in dynamic environments using density functions,''
arXiv:2411.12206, 2024.

\bibitem{morton2025oscbf}
D. Morton and M. Pavone,
``Safe, task-consistent manipulation with operational space
control barrier functions,''
in \emph{Proc. IEEE/RSJ Int. Conf. Intelligent Robots
and Systems (IROS)}, 2025, pp. 187--194.

\bibitem{xarmimplementation}
S. Krishnamoorthy,
``xarm-density: UF850 density-based control implementation.''
[Online]. Available:
\url{https://github.com/sriram-2502/xarm-density}

\end{thebibliography}

% =====================================================================
% Author checklist -- not printed
% =====================================================================
%
% Figures are loaded from figures/ with PDF preferred over PNG.
%
% Required basenames:
% fig01_tcp_tracking
% fig02_tracking_errors
% fig03_joint_velocities
% fig04_safety_intervention
% fig06_optimized_timing
% fig07_fixed_rate_budget
% fig08_robot_snapshots
%
% Tracking experiment: repository validation 9.
% Optimized benchmark and snapshots: repository validation 10.
%
% Short boxes are intentionally unbreakable.
% Tables are constrained to column width with tabularx.
% Float barriers keep all result figures before the conclusion
% and the obstacle table before the references.
%
% Before external submission:
% - replace the author placeholders;
% - record exact installed OSCBF and optimized-controller revisions;
% - verify the installed nominal-controller source matches the
%   explicit orientation and inverse definitions in this document;
% - update references to final publication metadata where appropriate.
%
% The reduction uses an orthonormal coordinate construction;
% an implementation may use an equivalent parameterization.

\end{document}
